% ====================================================================
% PRAKATA (PREFACE / FOREWORD)
% ====================================================================
% File ini berisi prakata/sekapur sirih (ucapan terima kasih).
% Prakata menjelaskan proses penelitian dan mengucapkan terima kasih kepada
% orang-orang yang telah membantu.
% Panjang: 0.5-1 halaman
% EDIT konten sesuai dengan penelitian Anda.
% ====================================================================

\frontmatterchapter{PRAKATA}

% --- ISI PRAKATA ---
\noindent Puji syukur ke hadirat Tuhan Yang Maha Esa atas segala berkah dan rahmat yang telah diberikan sehingga penelitian ini dapat diselesaikan dengan baik dan tepat pada waktunya.

Penelitian yang berjudul \textit{``\titleUpperCase''} ini merupakan salah satu persyaratan untuk memperoleh gelar Sarjana Komputer (S.Kom) pada Program Studi Informatika, Fakultas Ilmu Komputer, Universitas Jember.

Penulis menyadari bahwa dalam proses penelitian dan penulisan skripsi ini, banyak pihak yang telah memberikan bimbingan, dukungan, dan bantuan. Oleh karena itu, pada kesempatan ini penulis dengan sepenuh hati mengucapkan terima kasih kepada:

\begin{enumerate}
    \item \supervisor{} selaku Pembimbing Pertama yang telah memberikan bimbingan, arahan, dan motivasi
    \item Para dosen dan staf Program Studi Informatika yang telah membantu selama proses akademik
    \item Keluarga besar yang telah memberikan doa dan dukungan moral
    \item Rekan-rekan peneliti dan teman-teman yang telah memberikan masukan dan semangat
\end{enumerate}

Penulis menyadari bahwa skripsi ini masih memiliki berbagai kekurangan, baik dalam penulisan maupun isi. Oleh karena itu, penulis mengharapkan kritik dan saran yang sifatnya membangun dari pembaca untuk perbaikan di masa mendatang.

Akhir kata, semoga penelitian ini dapat memberikan manfaat bagi kemajuan ilmu pengetahuan dan dapat dijadikan sebagai bahan referensi bagi penelitian-penelitian selanjutnya.

\vspace{1cm}

Jember, \examinationDate

\vspace{2cm}

\noindent\authorName

\noindent\studentID

\newpage

% ====================================================================
% PETUNJUK PENGISIAN PRAKATA:
%
% 1. PANJANG: 0.5-1 halaman
%
% 2. STRUKTUR UMUM:
%    a) Ucapan syukur (1 paragraph)
%    b) Penjelasan judul skripsi dan tujuan (1 paragraph)
%    c) Daftar orang/institusi yang perlu diucapkan terima kasih
%       (dapat dalam bentuk paragraph atau numbered list)
%    d) Pernyataan kesadaran atas kekurangan
%    e) Harapan/doa penutup
%    f) Tempat dan Tanggal, Nama Penulis
%
% 3. ORANG YANG BIASANYA DISEBUT:
%    - Pembimbing akademik (I dan II)
%    - Dosen penguji
%    - Dosen lain yang memberikan kontribusi
%    - Institusi (universitas, departemen)
%    - Keluarga
%    - Teman-teman dekat/rekan peneliti
%
% 4. GAYA PENULISAN:
%    - Formal namun personal (tidak terlalu kaku)
%    - Ekspresif dan tulus
%    - Menggunakan bahasa yang hormat kepada yang ditebutkan
%
% 5. PENTING:
%    - Jangan lupa ucapan syukur di awal
%    - Pastikan nama pembimbing dan penguji spelled correctly
%    - Edit metadata untuk tanggal yang tepat
% ====================================================================
